% =====================================================
% 题型：解三角形角平分线最值解答题 (q_lyb_01_16_triangle_bisector.tex)
% =====================================================
% =====================================================
% 难度：★★★ (难题)
% =====================================================
\newcommand{\qTriangleBisectorMinStem}{%
  设 \(\triangle ABC\) 的内角 \(A, B, C\) 的对边分别为 \(a, b, c\)，\(\angle B\) 的角平分线 \(BD\) 与 \(AC\) 交于点 \(D\)。
  \begin{enumerate}[label=(\arabic*) , leftmargin=2em, itemsep=0.1em]
    \item 若 \(\vec{AD} = \dfrac{1}{5}\vec{AC}\)，求 \(\sin A : \sin C\) 的值；
    \item 若 \(b\cos C + \sqrt{3}b\sin C = a + c\)，且 \(BD = 2\)，求 \(c + 9a\) 的最小值。
  \end{enumerate}%
}

\newcommand{\qTriangleBisectorMinFigure}[1][1.3]{%
  \begin{tikzpicture}[scale=#1, >=stealth]
    % Coordinates of triangle
    \coordinate (B) at (0, 1.2);
    \coordinate (A) at (-1.0, 0);
    \coordinate (C) at (2.0, 0);
    
    % BD is the angle bisector.
    % D divides AC. Since AD = 1/5 AC => D is at 1/5 from A to C
    \coordinate (D) at ($(A)!0.2!(C)$);
    
    \draw[thick] (A) -- (B) -- (C) -- cycle;
    \draw[thick] (B) -- (D);
    
    \node[below left] at (A) {\small $A$};
    \node[above] at (B) {\small $B$};
    \node[below right] at (C) {\small $C$};
    \node[below] at (D) {\small $D$};
  \end{tikzpicture}%
}

\newcommand{\qTriangleBisectorMinAnswer}{%
  (1) \(\sin A : \sin C = 4:1\)；\\
  (2) \(c + 9a\) 的最小值为 \(\dfrac{32\sqrt{3}}{3}\)。%
}

\newcommand{\qTriangleBisectorMinSolution}{%
  (1) 因为 $\overrightarrow{AD} = \frac{1}{5}\overrightarrow{AC}$，所以 $AD = \frac{1}{5}AC$，
  
  从而有
  \[ AD : DC = 1 : 4. \]
  
  因为 $BD$ 是 $\angle B$ 的角平分线，由三角形角平分线定理得：
  \[ \frac{AB}{BC} = \frac{AD}{DC} = \frac{1}{4} \implies \frac{c}{a} = \frac{1}{4}. \]
  
  根据正弦定理，有：
  \[ \frac{\sin A}{\sin C} = \frac{a}{c} = 4. \]
  
  所以 $\sin A : \sin C$ 的值为 $4:1$。
  
  (2) 在 $\triangle ABC$ 中，由正弦定理将已知等式 $b\cos C + \sqrt{3}b\sin C = a + c$ 边化为角，得：
  \[ \sin B \cos C + \sqrt{3}\sin B \sin C = \sin A + \sin C. \]
  
  因为 $A = \pi - (B+C)$，所以
  \[ \sin A = \sin(B+C) = \sin B \cos C + \cos B \sin C. \]
  
  代入上式化简得：
  \[ \sin B \cos C + \sqrt{3}\sin B \sin C = \sin B \cos C + \cos B \sin C + \sin C, \]
  
  即
  \[ \sqrt{3}\sin B \sin C = \cos B \sin C + \sin C. \]
  
  因为 $C \in (0, \pi) \implies \sin C > 0$，等式两边同除以 $\sin C$，得：
  \[ \sqrt{3}\sin B - \cos B = 1, \]
  
  即
  \[ 2\sin\left(B - \frac{\pi}{6}\right) = 1 \implies \sin\left(B - \frac{\pi}{6}\right) = \frac{1}{2}. \]
  
  因为 $B \in (0, \pi) \implies B - \frac{\pi}{6} \in \left(-\frac{\pi}{6}, \frac{5\pi}{6}\right)$，
  
  所以有
  \[ B - \frac{\pi}{6} = \frac{\pi}{6} \implies B = \frac{\pi}{3}. \]
  
  因为 $BD = 2$ 是 $\angle B$ 的平分线，根据面积关系 $S_{\triangle ABC} = S_{\triangle ABD} + S_{\triangle CBD}$，有：
  \[ \frac{1}{2}ac \sin B = \frac{1}{2}c \cdot BD \sin\frac{B}{2} + \frac{1}{2}a \cdot BD \sin\frac{B}{2}. \]
  
  代入 $B = \frac{\pi}{3}$，$ \sin B = \frac{\sqrt{3}}{2}$，$ \sin\frac{B}{2} = \sin\frac{\pi}{6} = \frac{1}{2}$，以及 $BD = 2$，得：
  \[ \frac{\sqrt{3}}{4}ac = \frac{1}{2}c + \frac{1}{2}a \implies \frac{\sqrt{3}}{2}ac = a + c. \]
  
  两边同除以 $ac$，可得
  \[ \frac{1}{a} + \frac{1}{c} = \frac{\sqrt{3}}{2}. \]
  
  要求 $c + 9a$ 的最小值，利用常数代换法：
  \[
  \begin{aligned}
    c + 9a &= (c + 9a) \times \frac{2}{\sqrt{3}}\left(\frac{1}{a} + \frac{1}{c}\right) \\
    &= \frac{2}{\sqrt{3}} \left( \frac{c}{a} + 1 + 9 + \frac{9a}{c} \right) \\
    &= \frac{2}{\sqrt{3}}\left(10 + \frac{c}{a} + \frac{9a}{c}\right).
  \end{aligned}
  \]
  
  由均值不等式，有
  \[ \frac{c}{a} + \frac{9a}{c} \ge 2\sqrt{\frac{c}{a} \times \frac{9a}{c}} = 6, \]
  
  所以有
  \[ c + 9a \ge \frac{2}{\sqrt{3}}(10 + 6) = \frac{32}{\sqrt{3}} = \frac{32\sqrt{3}}{3}. \]
  
  当且仅当 $\frac{c}{a} = \frac{9a}{c} \implies c = 3a$ 时，等号成立。
  
  此时由 $\frac{1}{a} + \frac{1}{3a} = \frac{\sqrt{3}}{2} \implies \frac{4}{3a} = \frac{\sqrt{3}}{2} \implies a = \frac{8\sqrt{3}}{9}$，符合条件。
  
  所以 $c + 9a$ 的最小值为 $\frac{32\sqrt{3}}{3}$。
}

\newcommand{\qTriangleBisectorMinDiagnosis}{%
  学生第 (1) 问能用角平分线定理并得到正确比值；第 (2) 问中断，说明边角条件转化和最值链条尚未成型。关键缺口是由正弦定理推出 \(B=\frac{\pi}{3}\)，再由角平分线长度关系得到 \(\frac1a+\frac1c=\frac{\sqrt3}{2}\)，最后用均值不等式求最小值。%
}

\newcommand{\qTriangleBisectorMinTeachingNotes}{%
  精讲。课堂强调“边角化 -> 定角 -> 角平分线长度/面积关系 -> 倒数约束 -> 基本不等式”的完整链条。%
}


% =====================================================
% 别名兼容：旧课次宏定义
% =====================================================
\newcommand{\qLYBZeroOneSixteenStem}{\qTriangleBisectorMinStem}
\newcommand{\qLYBZeroOneSixteenFigure}[1][1.2]{\qTriangleBisectorMinFigure[#1]}
\newcommand{\qLYBZeroOneSixteenAnswer}{\qTriangleBisectorMinAnswer}
\newcommand{\qLYBZeroOneSixteenSolution}{\qTriangleBisectorMinSolution}
\newcommand{\qLYBZeroOneSixteenDiagnosis}{\qTriangleBisectorMinDiagnosis}
\newcommand{\qLYBZeroOneSixteenTeachingNotes}{\qTriangleBisectorMinTeachingNotes}
